% Pursuing Stacks (cotes 134-2 to 134-8): archive cote 134-7, site page 64 (PDF page 65), typescript page 506.
% typescript. Draft compared: ps6.tex:3120-3190.
% AI-transcribed by Claude Opus 5.5 (claude-opus-5-5, adaptive thinking, effort high), Message Batches API, batch pilot (msgbatch_0189Nhf1xB615YWnztubtLfc), 2026-11-12; not yet checked by eye.
\scanpage{134-7/64}
\mspage{506}

28.9.
\inserted{(prompted by the reflections from the day before, cf. section 119)}
120.\quad Yesterday I pondered a little on the common features of the various
set-ups for ``commutative algebra'' (possibly too for corresponding notions
of ``schemes'') one gets when introducing extra operations on commutative
rings or algebras, such as divided power structure\xout{s} (on a suitable ideal\hand{),}
or a \hand{$\lambda$}-structure with\xout{\unsure{n}} operations \hand{$\lambda$}$^i$ paraphrasing exterior powers, or a
binomial structure with operations $x\hand{$\mapsto$}\binom{x}{n}$ paraphrasing binomial coeffi-%
cients, or an S-structure with operations $S^i$ of Adam's type paraphrasing
sums of i.th powers of roots of a polynomial (a weakened version of a \hand{$\lambda$}-%
structure). It seems that the unifying notion here is the notion of an
``analyzer'' (analyseur) of Lazare, ``containing'' the Lazare analyzer for
commutative rings (not necessarily with unit), so that the components
\margin{$n\geqslant -1$} \hand{$\Omega$}$_n$ (n\hand{$\geqslant$}\over{1}{-1}) are commutative rings (not necessarily with unit). In case
the \xout{\unsure{structural}}\inserted[typed]{extra} operations we want to introduce on commutative rings are
to be defined on rings \uline{with units} (not just on a suitable ideal of such
a ring, as is the case for the divided power structure), and they all
can be defined in terms of operations involving just one argument, the
reasonable extra axiom on \xout{the} the corresponding analyzer (as suggested by
the examples at hand) is that for n \hand{$\geqslant$} \over{\unsure{2}}{1}, \hand{$\Omega$}$_n$ can be recovered in terms
of \hand{$\Omega$}$_0$ and \hand{$\Omega$}$_{-1}$ = k$_0$ (the latter act\xout{\unsure{ing}}\inserted[typed]{ing} as a ground ring for the theory\hand{)}
as the \deleted{n}\inserted{(n+1)}-fold tensor power of \hand{$\Omega$}$_0$ over k$_0$ . Thus the \deleted{\illegible{1}}whole structure of
\margin{à la ligne} the analyzor may be thought of as embodied in the system \inserted{$\Omega = (k_0, \Omega_0)$,} \deleted{$(k_0, \Omega_0)$, where}
\inserted{where} k$_0$ is a commutative ring (with unit, now), \hand{$\Omega$}$_0$ a commutative k$_0$-algebra
\margin{id} with unit, endowed moreover with a composition operation \deleted{$(F,G)\mapsto F\circ G$,}\inserted{$(F,G)\longmapsto F\circ G$ ,}
satisfying a bunch of simple axioms I do'nt feel like writing down here.
The simplest case of all of course (correspo\unsure{n}ding to \xout{the} \xout{u}usual commutati%
ve algebra ``over k$_0$'' as a ground ring, with no extra structure on commu-%
tative algebras with unit over k$_0$) is \hand{$\Omega$}$_0$ = k\hand{[}T\hand{]}\hand{,} with the usual compositi%
on of polynomials, T acting as the two-sided unit for composition. In
the general case, \hand{$\Omega$}$_0$ and it's tensor powers \inserted{$\Omega_n$ over $k_0$} are going to play the part
played by polynomial rings in ordinary commutative algebra, \deleted{and} \over{t}{T}here
should be a ready generalization, in this spirit, of taking the symme-%
tric algebra of a k$_0$-module (which, for a module free and of finite type\hand{\unsure{,}}
will \deleted{be}\inserted{yield a $\Omega$-structure} isomorphic to one of the \hand{$\Omega$}$_n$'s). An \hand{$\Omega$}-structure on a set \xout{x} k
amounts to giving a structure of a commutative k$_0$-algebra with unit on k,
plus a map
\[ \hand{$\Omega_0 \to$} \mathrm{Maps}(k,k) \quad \hand{$,$} \quad \hand{$F \longmapsto (x \longmapsto F(x))$} \]

\margin{64}
